\section{Background}\label{sec:bg}

Fix a network of timed automata with processes $P_1,\dots,P_n$ over clocks $\clockSet$,
$|\clockSet| = m$, and bounded integer variables. A \emph{discrete configuration} $\ell$
consists of one location per process together with the integer valuation; $\mathcal{S}$
denotes the set of reachable configurations. A \emph{zone} is a set of clock valuations
defined by difference constraints $x_i - x_j \le c$, and a \emph{symbolic state} is a pair
$(\ell,\Zone)$, standing for infinitely many concrete states at once; $\Sigma$ denotes the set
of reachable symbolic states. Reachability is decided by exploring the \emph{zone graph} whose
nodes are symbolic states.

A zone over $m$ clocks is stored as a Difference Bound Matrix: an $(m{+}1) \times (m{+}1)$
matrix $\DBM{}$ whose entry $(i,j)$ is the tightest known bound on $x_i - x_j$, with index $0$
representing the constant zero. A DBM is kept \emph{canonical}, meaning every entry is already
the tightest implied bound, which makes inclusion testing a pointwise comparison of matrices.
Canonicalisation is all-pairs shortest paths, $O(m^3)$ by Floyd--Warshall, but adding a single
constraint can be re-canonicalised incrementally in $O(m^2)$; a zone is empty exactly when
some diagonal entry is negative. Computing the successor of $(\ell,\Zone)$ along an edge
intersects the guard, applies clock resets, lets time elapse, intersects the target
invariant, and finally applies the $\aLU{}$ extrapolation that keeps the zone graph finite.
All of these are $O(m^2)$ except the emptiness test, which is $O(m)$.

Exploration follows the classical scheme with a \emph{visited} set of expanded states and a
\emph{waiting} set of discovered-but-unexpanded ones. We proceed level-synchronously, so the
waiting set is simply the \emph{frontier}: the symbolic states discovered during the previous
breadth-first level, all of which are expanded during the current one. One level therefore
expands every state in the frontier, inserts the resulting successors into the visited set,
and collects those that survive insertion into the frontier of the next level; levels are
separated by a barrier.

Insertion is where timed model checking parts company with the untimed case. A new
$(\ell,\Zone)$ is redundant, and discarded, if some stored $(\ell,\Zone')$ has
$\Zone \subseteq \Zone'$; symmetrically it may render stored zones redundant, and they are
removed. In place of plain inclusion we use the $\aLU{}$ abstraction, parameterised by
per-clock bounds derived from $\ell$. Additionally, when the union of two zones happens to be
convex it is representable exactly, and the pair is replaced by that union. What is stored for
a configuration is therefore an \emph{antichain} of mutually incomparable zones, and its
length is precisely the cost of inserting into it.

This yields the one structural quantity used throughout the paper,
\[
  \rho \;=\; |\Sigma| \,/\, |\mathcal{S}|,
\]
the mean number of zones per discrete configuration, equivalently the mean antichain length.
Benchmarks differ enormously along this axis, and it will turn out to decide not only how
expensive insertion is but which of our two engines should be used. Viking sits at
$\rho \approx 1$, with essentially one zone per configuration, whereas Fischer concentrates
many zones on comparatively few configurations. Prior work already identified $\rho$ as
governing the scalability of a shared visited set; \cref{sec:batching} shows how to stop it
from doing so.
